\documentclass[UTF8, 12pt, fontset=fandol, oneside]{ctexart}
\usepackage{researchnotes}

\title{常见初等不等式}
\author{}
\date{}

\begin{document}

\maketitle

不等式是初等数学与数学分析中最常用的估计工具. 在许多问题中，我们并不需要精确求值，而只需要知道一个量不超过什么、至少有多大、何时取到最优值. 因此，学习不等式的关键不仅是记住结论，还要理解它们分别适合处理哪一种结构：加法、乘法、平方和、平均、排序、凸性、指数对数或三角函数.

本文整理若干常见初等不等式，并附上适用条件、等号条件和典型用法. 其中默认变量均为实数；若出现几何平均、调和平均、对数、分母等表达式，则需额外满足正性或非零条件.

\section*{一、不等式的基本原则}

\noindent\textbf{原则 1（同向相加）}\quad 若
\[
  a\le b,\qquad c\le d,
\]
则
\[
  a+c\le b+d.
\]
这说明多个上界估计可以逐项相加.

\noindent\textbf{原则 2（乘以正数）}\quad 若 $a\le b$，$\lambda>0$，则
\[
  \lambda a\le \lambda b.
\]
若 $\lambda<0$，则不等号方向反向：
\[
  \lambda a\ge \lambda b.
\]
很多错误都来自忘记负数乘法会改变不等号方向.

\noindent\textbf{原则 3（平方非负）}\quad 对任意实数 $x$，
\[
  x^2\ge0.
\]
更一般地，
\[
  (a-b)^2\ge0.
\]
这是大量初等不等式的源头. 例如
\[
  a^2+b^2\ge2ab,
\]
等号当且仅当 $a=b$.

\noindent\textbf{原则 4（单调函数保序）}\quad 若函数 $f$ 在区间 $I$ 上单调递增，且 $a,b\in I,\ a\le b$，则
\[
  f(a)\le f(b).
\]
若 $f$ 单调递减，则
\[
  f(a)\ge f(b).
\]
例如，指数函数 $e^x$ 单调递增，对数函数 $\ln x$ 在 $(0,+\infty)$ 上单调递增，而倒数函数 $1/x$ 在 $(0,+\infty)$ 上单调递减.

\section*{二、绝对值与三角不等式}

\noindent\textbf{定义}\quad 实数 $x$ 的绝对值为
\[
  |x|=
  \begin{cases}
    x, & x\ge0,\\
    -x, & x<0.
  \end{cases}
\]
因此
\[
  -|x|\le x\le |x|.
\]

\noindent\textbf{三角不等式}\quad 对任意实数 $x,y$，
\[
  |x+y|\le |x|+|y|.
\]
更一般地，
\[
  \left|\sum_{i=1}^n x_i\right|
  \le
  \sum_{i=1}^n |x_i|.
\]

\noindent\textbf{反三角不等式}\quad 对任意实数 $x,y$，
\[
  \bigl||x|-|y|\bigr|\le |x-y|.
\]
这个不等式常用于估计绝对值函数的变化量.

\noindent\textbf{常用推论}\quad 若 $a>0$，则
\[
  |x|\le a
  \quad\Longleftrightarrow\quad
  -a\le x\le a.
\]
若 $|x-a|\le r$，则 $x$ 落在以 $a$ 为中心、半径为 $r$ 的闭区间中：
\[
  a-r\le x\le a+r.
\]

\section*{三、平方型不等式}

\noindent\textbf{基本平方不等式}\quad 对任意实数 $a,b$，
\[
  a^2+b^2\ge2ab.
\]
等号当且仅当 $a=b$.

将 $b$ 换成 $-b$，可得
\[
  a^2+b^2\ge -2ab.
\]
因此
\[
  2|ab|\le a^2+b^2.
\]

\noindent\textbf{带参数形式}\quad 对任意实数 $a,b$ 和任意 $\varepsilon>0$，
\[
  2|ab|\le \varepsilon a^2+\frac{1}{\varepsilon}b^2.
\]
证明只需注意
\[
  \left(\sqrt{\varepsilon}|a|-\frac{|b|}{\sqrt{\varepsilon}}\right)^2\ge0.
\]
这个形式常用于把乘积项拆成两个平方项，并且可以通过选择 $\varepsilon$ 调整两部分权重.

\noindent\textbf{平方和估计}\quad 对任意实数 $x_1,x_2,\cdots,x_n$，
\[
  \left(\sum_{i=1}^n x_i\right)^2
  \le
  n\sum_{i=1}^n x_i^2.
\]
这是 Cauchy--Schwarz 不等式的一个重要特例. 当 $x_1=x_2=\cdots=x_n$ 时取等号.

\section*{四、平均值不等式}

\noindent\textbf{算术--几何平均不等式（AM--GM）}\quad 若 $x_1,x_2,\cdots,x_n>0$，则
\[
  \frac{x_1+x_2+\cdots+x_n}{n}
  \ge
  \sqrt[n]{x_1x_2\cdots x_n}.
\]
等号当且仅当
\[
  x_1=x_2=\cdots=x_n.
\]

\noindent\textbf{二元形式}\quad 若 $a,b\ge0$，则
\[
  \frac{a+b}{2}\ge\sqrt{ab},
\]
等价于
\[
  a+b\ge2\sqrt{ab}.
\]
这是由 $(\sqrt a-\sqrt b)^2\ge0$ 直接推出的.

\noindent\textbf{加权 AM--GM}\quad 若 $x_i>0$，$w_i\ge0$，且
\[
  \sum_{i=1}^n w_i=1,
\]
则
\[
  \sum_{i=1}^n w_i x_i
  \ge
  \prod_{i=1}^n x_i^{w_i}.
\]
当所有正权重对应的 $x_i$ 相等时取等号.

\noindent\textbf{调和--几何--算术--平方平均不等式}\quad 对正数 $x_1,x_2,\cdots,x_n$，定义
\[
  H=\frac{n}{\frac1{x_1}+\frac1{x_2}+\cdots+\frac1{x_n}},
  \qquad
  G=\sqrt[n]{x_1x_2\cdots x_n},
\]
\[
  A=\frac{x_1+x_2+\cdots+x_n}{n},
  \qquad
  Q=\sqrt{\frac{x_1^2+x_2^2+\cdots+x_n^2}{n}}.
\]
则
\[
  H\le G\le A\le Q.
\]
等号同时成立当且仅当
\[
  x_1=x_2=\cdots=x_n.
\]

\noindent\textbf{幂平均不等式}\quad 对正数 $x_1,\cdots,x_n$，定义
\[
  M_r=
  \left(\frac{x_1^r+x_2^r+\cdots+x_n^r}{n}\right)^{1/r}
  \quad(r\ne0),
\]
并定义
\[
  M_0=\sqrt[n]{x_1x_2\cdots x_n}.
\]
若 $r<s$，则
\[
  M_r\le M_s.
\]
上述 $H,G,A,Q$ 分别对应 $r=-1,0,1,2$ 的幂平均.

\section*{五、Cauchy--Schwarz 不等式}

\noindent\textbf{定理（Cauchy--Schwarz）}\quad 对任意实数 $a_1,\cdots,a_n$ 和 $b_1,\cdots,b_n$，
\[
  \left(\sum_{i=1}^n a_i b_i\right)^2
  \le
  \left(\sum_{i=1}^n a_i^2\right)
  \left(\sum_{i=1}^n b_i^2\right).
\]
等号成立当且仅当两个向量成比例，即存在常数 $\lambda$，使得
\[
  a_i=\lambda b_i\qquad (i=1,2,\cdots,n),
\]
其中零向量情形需单独理解为显然取等号.

\noindent\textbf{常用特例}\quad 取 $b_1=b_2=\cdots=b_n=1$，得
\[
  \left(\sum_{i=1}^n a_i\right)^2
  \le
  n\sum_{i=1}^n a_i^2.
\]

\noindent\textbf{Engel 形式}\quad 若 $x_i\ge0$，$a_i>0$，则
\[
  \sum_{i=1}^n \frac{x_i^2}{a_i}
  \ge
  \frac{(x_1+x_2+\cdots+x_n)^2}{a_1+a_2+\cdots+a_n}.
\]
特别地，若 $x_i>0$，则
\[
  \sum_{i=1}^n \frac1{x_i}
  \ge
  \frac{n^2}{x_1+x_2+\cdots+x_n}.
\]
这正是调和平均不超过算术平均的一个证明.

\noindent\textbf{例}\quad 若 $a,b,c>0$，则
\[
  \frac{a^2}{b+c}+\frac{b^2}{c+a}+\frac{c^2}{a+b}
  \ge
  \frac{(a+b+c)^2}{2(a+b+c)}
  =
  \frac{a+b+c}{2}.
\]
这类分式求和估计通常优先尝试 Engel 形式.

\section*{六、Holder 与 Minkowski 不等式}

\noindent\textbf{Holder 不等式}\quad 设 $p,q>1$，且
\[
  \frac1p+\frac1q=1.
\]
若 $a_i,b_i\ge0$，则
\[
  \sum_{i=1}^n a_i b_i
  \le
  \left(\sum_{i=1}^n a_i^p\right)^{1/p}
  \left(\sum_{i=1}^n b_i^q\right)^{1/q}.
\]
当 $p=q=2$ 时，它就是 Cauchy--Schwarz 不等式.

\noindent\textbf{三列常用形式}\quad 若 $a_i,b_i,c_i\ge0$，则
\[
  \left(\sum_{i=1}^n a_i b_i c_i\right)^3
  \le
  \left(\sum_{i=1}^n a_i^3\right)
  \left(\sum_{i=1}^n b_i^3\right)
  \left(\sum_{i=1}^n c_i^3\right).
\]

\noindent\textbf{Minkowski 不等式}\quad 设 $p\ge1$，则
\[
  \left(\sum_{i=1}^n |a_i+b_i|^p\right)^{1/p}
  \le
  \left(\sum_{i=1}^n |a_i|^p\right)^{1/p}
  +
  \left(\sum_{i=1}^n |b_i|^p\right)^{1/p}.
\]
当 $p=1$ 时，它就是三角不等式的求和形式；当 $p=2$ 时，它是欧氏距离中的三角不等式.

\section*{七、Young 不等式}

\noindent\textbf{定理（Young 不等式）}\quad 设 $a,b\ge0$，$p,q>1$，且
\[
  \frac1p+\frac1q=1.
\]
则
\[
  ab\le \frac{a^p}{p}+\frac{b^q}{q}.
\]
等号成立当且仅当
\[
  a^p=b^q.
\]

\noindent\textbf{常用特例}\quad 当 $p=q=2$ 时，
\[
  ab\le \frac{a^2+b^2}{2}.
\]
这就是平方型不等式 $a^2+b^2\ge2ab$ 的另一种写法.

\noindent\textbf{带参数形式}\quad 对任意 $a,b\in\mathbb R$ 和任意 $\varepsilon>0$，
\[
  |ab|\le \frac{\varepsilon a^2}{2}+\frac{b^2}{2\varepsilon}.
\]
它常用于分析中吸收交叉项，例如把 $|ab|$ 控制到 $a^2$ 与 $b^2$ 的组合中.

\section*{八、Bernoulli 不等式}

\noindent\textbf{定理（Bernoulli 不等式）}\quad 若 $x\ge -1$，$n$ 为正整数，则
\[
  (1+x)^n\ge1+nx.
\]
当 $n=1$ 时恒取等号；当 $n>1$ 时，等号当且仅当 $x=0$.

\noindent\textbf{证明思路}\quad 用数学归纳法. 若
\[
  (1+x)^n\ge1+nx,
\]
则因为 $1+x\ge0$，
\[
\begin{aligned}
  (1+x)^{n+1}
  &\ge (1+nx)(1+x)\\
  &=1+(n+1)x+nx^2\\
  &\ge1+(n+1)x.
\end{aligned}
\]

\noindent\textbf{实指数形式}\quad 若 $x>-1$，则
\[
  (1+x)^\alpha\ge1+\alpha x
\]
在 $\alpha\ge1$ 或 $\alpha\le0$ 时成立；而当 $0\le\alpha\le1$ 时，不等号方向相反：
\[
  (1+x)^\alpha\le1+\alpha x.
\]
这可以由函数 $(1+x)^\alpha$ 的凸凹性得到.

\section*{九、指数与对数不等式}

\noindent\textbf{指数函数的切线不等式}\quad 对任意实数 $x$，
\[
  e^x\ge1+x.
\]
等号当且仅当 $x=0$.

更一般地，由 Taylor 展开或凸性可得
\[
  e^x\ge1+x+\frac{x^2}{2}\qquad(x\ge0).
\]

\noindent\textbf{对数函数的切线不等式}\quad 对任意 $x>0$，
\[
  \ln x\le x-1.
\]
等号当且仅当 $x=1$.

令 $x=1+t$，得到常用形式
\[
  \ln(1+t)\le t,\qquad t>-1.
\]

\noindent\textbf{对数下界}\quad 当 $t\ge0$ 时，
\[
  \ln(1+t)\ge \frac{t}{1+t}.
\]
证明可设
\[
  f(t)=\ln(1+t)-\frac{t}{1+t},
\]
则
\[
  f'(t)=\frac{t}{(1+t)^2}\ge0,\qquad f(0)=0.
\]

\noindent\textbf{常用夹估计}\quad 当 $t\ge0$ 时，
\[
  \frac{t}{1+t}\le \ln(1+t)\le t.
\]
它常用于估计含对数的极限与级数.

\section*{十、排序型不等式}

\noindent\textbf{重排不等式}\quad 设
\[
  a_1\le a_2\le\cdots\le a_n,\qquad
  b_1\le b_2\le\cdots\le b_n.
\]
则在所有排列 $\sigma$ 中，
\[
  \sum_{i=1}^n a_i b_{n+1-i}
  \le
  \sum_{i=1}^n a_i b_{\sigma(i)}
  \le
  \sum_{i=1}^n a_i b_i.
\]
即同序相乘的和最大，反序相乘的和最小.

\noindent\textbf{Chebyshev 求和不等式}\quad 若
\[
  a_1\le a_2\le\cdots\le a_n,\qquad
  b_1\le b_2\le\cdots\le b_n,
\]
则
\[
  \frac1n\sum_{i=1}^n a_i b_i
  \ge
  \left(\frac1n\sum_{i=1}^n a_i\right)
  \left(\frac1n\sum_{i=1}^n b_i\right).
\]
如果一列递增而另一列递减，则不等号方向反向.

\noindent\textbf{适用场景}\quad 排序型不等式适合处理“大的配大的，小的配小的”这类结构. 在应用前，必须先确认两组数具有相同或相反的排序方向.

\section*{十一、Jensen 不等式与凸性方法}

\noindent\textbf{凸函数定义}\quad 设 $f$ 是区间 $I$ 上的函数. 若对任意 $x,y\in I$ 和 $\lambda\in[0,1]$，
\[
  f(\lambda x+(1-\lambda)y)
  \le
  \lambda f(x)+(1-\lambda)f(y),
\]
则称 $f$ 为凸函数.

\noindent\textbf{Jensen 不等式}\quad 若 $f$ 是区间 $I$ 上的凸函数，$x_i\in I$，$w_i\ge0$，且
\[
  \sum_{i=1}^n w_i=1,
\]
则
\[
  f\left(\sum_{i=1}^n w_i x_i\right)
  \le
  \sum_{i=1}^n w_i f(x_i).
\]
若 $f$ 为凹函数，则不等号方向反向.

\noindent\textbf{典型例子}\quad 因为 $f(x)=x^2$ 是凸函数，所以
\[
  \left(\frac{x_1+\cdots+x_n}{n}\right)^2
  \le
  \frac{x_1^2+\cdots+x_n^2}{n}.
\]
因为 $f(x)=\ln x$ 在 $(0,+\infty)$ 上是凹函数，所以
\[
  \ln\left(\frac{x_1+\cdots+x_n}{n}\right)
  \ge
  \frac{\ln x_1+\cdots+\ln x_n}{n},
\]
从而得到 AM--GM：
\[
  \frac{x_1+\cdots+x_n}{n}
  \ge
  \sqrt[n]{x_1x_2\cdots x_n}.
\]

\section*{十二、三角函数不等式}

\noindent\textbf{基本三角估计}\quad 当 $0<x<\dfrac{\pi}{2}$ 时，
\[
  \sin x<x<\tan x.
\]
等价地，
\[
  \cos x<\frac{\sin x}{x}<1.
\]
令 $x\to0^+$，可以得到重要极限
\[
  \lim_{x\to0}\frac{\sin x}{x}=1.
\]

\noindent\textbf{全局估计}\quad 对任意实数 $x$，
\[
  |\sin x|\le |x|.
\]
这是由 $\sin x$ 的导数绝对值不超过 $1$ 或由基本三角估计推广得到的.

\noindent\textbf{余弦估计}\quad 对任意实数 $x$，
\[
  1-\frac{x^2}{2}\le \cos x\le1.
\]
其中右边来自 $|\cos x|\le1$ 的一部分；左边可由 Taylor 展开余项或函数
\[
  f(x)=\cos x-1+\frac{x^2}{2}
\]
的性质证明.

\noindent\textbf{常用小量估计}\quad 当 $x$ 接近 $0$ 时，
\[
  \sin x\sim x,\qquad 1-\cos x\sim \frac{x^2}{2}.
\]
若只需不等式估计，可使用
\[
  0\le1-\cos x\le\frac{x^2}{2}.
\]

\section*{十三、常用组合与例题}

\noindent\textbf{例 1}\quad 证明若 $x>0$，则
\[
  x+\frac1x\ge2.
\]

\noindent\textbf{证明}\quad 由 AM--GM，
\[
  \frac{x+\frac1x}{2}\ge \sqrt{x\cdot\frac1x}=1.
\]
故
\[
  x+\frac1x\ge2.
\]
等号当且仅当 $x=1$.

\noindent\textbf{例 2}\quad 证明若 $a,b,c>0$，则
\[
  (a+b+c)\left(\frac1a+\frac1b+\frac1c\right)\ge9.
\]

\noindent\textbf{证明}\quad 由 Cauchy--Schwarz 不等式，
\[
  (a+b+c)\left(\frac1a+\frac1b+\frac1c\right)
  \ge
  (1+1+1)^2=9.
\]
等号当且仅当 $a=b=c$.

\noindent\textbf{例 3}\quad 证明若 $x_i\ge0$，则
\[
  \prod_{i=1}^n (1+x_i)
  \ge
  1+\sum_{i=1}^n x_i.
\]

\noindent\textbf{证明}\quad 展开左边：
\[
  \prod_{i=1}^n(1+x_i)
  =
  1+\sum_{i=1}^n x_i+\text{若干个非负的高阶乘积项}.
\]
因此结论成立. 等号成立当且仅当任意两个 $x_i$ 不同时为正.

\noindent\textbf{例 4}\quad 证明若 $x>-1$，则
\[
  \ln(1+x)\le x.
\]

\noindent\textbf{证明}\quad 由对数切线不等式 $\ln y\le y-1$，令 $y=1+x>0$，得
\[
  \ln(1+x)\le x.
\]

\section*{十四、选用不等式的思路}

\noindent\textbf{含绝对值}\quad 优先考虑三角不等式、反三角不等式以及 $|ab|\le\dfrac{a^2+b^2}{2}$.

\noindent\textbf{含平方和}\quad 优先考虑 Cauchy--Schwarz、不等式
\[
  \left(\sum x_i\right)^2\le n\sum x_i^2
\]
以及平方非负.

\noindent\textbf{含正数乘积}\quad 优先考虑 AM--GM、加权 AM--GM、Bernoulli 不等式.

\noindent\textbf{含分式求和}\quad 优先考虑 Engel 形式：
\[
  \sum \frac{x_i^2}{a_i}\ge\frac{(\sum x_i)^2}{\sum a_i}.
\]

\noindent\textbf{含指数对数}\quad 优先考虑
\[
  e^x\ge1+x,\qquad \ln x\le x-1,\qquad \ln(1+x)\le x.
\]

\noindent\textbf{含平均或凸函数}\quad 优先考虑 Jensen 不等式. 若函数二阶导数非负，通常使用凸函数版本；若二阶导数非正，则使用凹函数版本.

\section*{十五、小结}

常见初等不等式可以按结构记忆：
\[
\begin{array}{c|c}
  \text{结构} & \text{常用不等式}\\ \hline
  \text{绝对值} & \text{三角不等式、反三角不等式}\\
  \text{平方} & (a-b)^2\ge0,\ a^2+b^2\ge2ab\\
  \text{平均} & H\le G\le A\le Q,\ \text{幂平均}\\
  \text{内积与平方和} & \text{Cauchy--Schwarz}\\
  \text{乘积与幂} & \text{AM--GM、Young、Bernoulli}\\
  \text{指数对数} & e^x\ge1+x,\ \ln x\le x-1\\
  \text{排序} & \text{重排不等式、Chebyshev 不等式}\\
  \text{凸性} & \text{Jensen 不等式}\\
  \text{三角函数} & \sin x<x<\tan x,\ |\sin x|\le|x|
\end{array}
\]
使用不等式时，最重要的是先识别表达式的结构，再检查适用条件和等号条件. 一个好的不等式证明，往往不是把结论强行变形，而是找到它背后的“平均”“平方”“凸性”或“排序”结构.

\end{document}
